% 宏观 QED：有损介质的连续储库量子化

Maxwell 方程把线性材料响应压缩到电磁算符 \(\mathsf L_{\rm EM}\)、推迟并矢 Green 张量 \(\mathbf G^R\) 与相应光学恒等式。存在吸收时，电磁子系统本身不是封闭哈密顿系统；量子化必须显式加入环境储库，使场算符既保持正则对易关系，又让真空与热涨落由同一个耗散核 \(\operatorname{Im}\mathbf G^R\) 决定。

对只含局域电损耗的介质，引入连续玻色储库算符
\begin{equation*}
 \hat{\mathbf f}(\mathbf r,\omega),
 \qquad \omega>0,
\end{equation*}
并规定
\begin{equation}
 \boxed{
 [\hat f_i(\mathbf r,\omega),
 \hat f_j^\dagger(\mathbf r',\omega')]
 =\delta_{ij}\delta^{(3)}(\mathbf r-\mathbf r')\delta(\omega-\omega').}
 \label{eq:reservoir-comm-SI-r12}
\end{equation}
介质的吸收谱进入噪声电流
\begin{equation}
 \boxed{
 \hat{\mathbf J}_{N}(\mathbf r,\omega)
 =\omega
 \sqrt{\frac{\hbar\varepsilon_0}{\pi}
 \operatorname{Im}\boldsymbol\varepsilon_r(\mathbf r,\omega)}
 \cdot\hat{\mathbf f}(\mathbf r,\omega),}
 \label{eq:noise-current-SI-r12}
\end{equation}
其中矩阵平方根取半正定分支。\eqnrefs{eq:noise-current-SI-r12} 的归一化由电场正则对易关系与 Green 张量光学恒等式共同固定，而不是独立的唯象噪声强度。


把噪声电流代入 Maxwell 推迟解，得到正频电场算符
\begin{equation}
 \boxed{
 \hat{\mathbf E}^{(+)}(\mathbf r,\omega)
 =\ii\mu_0\omega^2
 \int\dd^3 s\,
 \mathbf G^R(\mathbf r,\mathbf s;\omega)
 \cdot
 \sqrt{\frac{\hbar\varepsilon_0}{\pi}
 \operatorname{Im}\boldsymbol\varepsilon_r(\mathbf s,\omega)}
 \cdot\hat{\mathbf f}(\mathbf s,\omega).}
 \label{eq:macqed-Eplus-SI-r12}
\end{equation}
完整时域场由正、负频部分叠加：
\begin{equation*}
 \hat{\mathbf E}(\mathbf r,t)
 =\int_0^\infty\dd\omega\,
 \hat{\mathbf E}^{(+)}(\mathbf r,\omega)\ee^{-\ii\omega t}
 +\mathrm{H.c.}
\end{equation*}
因此介质的经典吸收谱不仅决定平均响应，也决定量子场涨落能够从哪些空间与频率通道进入。

有限温度 \(T_{\rm th}>0\) 下，定义 Bose--Einstein 占据数
\begin{equation*}
 n_B(\omega;T_{\rm th})
 =\frac1{\exp[\hbar\omega/(k_BT_{\rm th})]-1}.
\end{equation*}
先把热态中每个模式的平均占据写成
\begin{equation}
\boxed{
 \langle\hat f_i^\dagger(\mathbf r,\omega)
 \hat f_j(\mathbf r',\omega')\rangle_{T_{\rm th}}
 =n_B(\omega;T_{\rm th})\,
 \delta_{ij}\delta^{(3)}(\mathbf r-\mathbf r')\delta(\omega-\omega').}
\label{eq:reservoir-thermal-occupation-dp68}
\end{equation}
再结合玻色对易关系，热平衡储库的两种有序相关函数为
\begin{align}
 \langle\hat f_i^\dagger(\mathbf r,\omega)
 \hat f_j(\mathbf r',\omega')\rangle_{T_{\rm th}}
 &=n_B(\omega;T_{\rm th})\,
 \delta_{ij}\delta^{(3)}(\mathbf r-\mathbf r')\delta(\omega-\omega'),\notag\\
 \langle\hat f_i(\mathbf r,\omega)
 \hat f_j^\dagger(\mathbf r',\omega')\rangle_{T_{\rm th}}
 &=[n_B(\omega;T_{\rm th})+1]\,
 \delta_{ij}\delta^{(3)}(\mathbf r-\mathbf r')\delta(\omega-\omega').
 \label{eq:reservoir-thermal-correlators-dp43}
\end{align}
第二行比第一行多出的单位项来自玻色对易关系；它区分环境吸收量子与向系统发射量子的两种有序相关函数。零温时 \(n_B\to0\)，储库仍保留这个单位项。把它代入\eqnrefs{eq:macqed-Eplus-SI-r12}，再用 Green 张量光学恒等式，可得
\begin{equation}
\boxed{
 \langle0|\hat{\mathbf E}^{(+)}(\mathbf r,\omega)
 \hat{\mathbf E}^{(-)}(\mathbf r',\omega')|0\rangle
 =\frac{\hbar\mu_0}{\pi}\omega^2
 \operatorname{Im}\mathbf G^R(\mathbf r,\mathbf r';\omega)
 \delta(\omega-\omega').}
\label{eq:macqed-vacuum-field-correlation-dp68}
\end{equation}
因此同一个 \(\operatorname{Im}\mathbf G^R\) 一方面描述经典吸收，另一方面固定零点电场涨落的谱权重。



若系统无损、封闭，并且 Maxwell 广义本征问题能够化为厄米形式，则推迟 Green 张量退化为离散实频率极点之和：
\begin{equation}
 \mathbf G^R(\mathbf r,\mathbf r';\omega)
 =\sum_\lambda
 \frac{\mathbf E_\lambda(\mathbf r)\otimes\mathbf E_\lambda^*(\mathbf r')}
 {\mathcal N_\lambda[\omega_\lambda^2-(\omega+\ii0^+)^2]}.
 \label{eq:G-mode-expansion-SI-r12}
\end{equation}
这里 \(\mathcal N_\lambda\) 定义为相应极点留数的电磁能量归一化。色散介质中它一般包含
\begin{equation*}
 \frac{\partial[\omega\epsilon(\omega)]}{\partial\omega},
 \qquad
 \frac{\partial[\omega\mu(\omega)]}{\partial\omega},
\end{equation*}
不能直接照搬真空中的 \(\int(\epsilon|\mathbf E|^2+\mu|\mathbf H|^2)\)。一旦存在真实吸收，离散正常模展开不再完备，必须回到连续储库与 Green 张量表示。


有损介质比无损正常模体系多出一个不可缺少的对象：连续储库。它与经典 Maxwell Green 张量并不是两套独立理论；材料吸收先进入 \(\operatorname{Im}\boldsymbol\varepsilon_r\)，随后通过\eqnrefs{eq:green-optical-identity-SI-r12} 决定 \(\operatorname{Im}\mathbf G^R\)，最终由\eqnrefs{eq:macqed-vacuum-field-correlation-dp68} 固定量子场涨落。无损封闭系统的离散正常模只是在耗散消失后的特殊极限。

